% Primene dinamickih grafova

Динамички графови налазе широку примену у системима где се топологија мреже континуирано мења кроз време. Са порастом мрежне инфраструктуре, развојем науке и технологије потреба за ефикасним ажурирањима и одржавањем стања различитих мрежа у реалном времену је све више заступљена. 

У наредним секцијама биће представљене само неке од бројних примена динамичких графова како у науци тако и у индустрији, што додатно даје значај овој теми.

% Dinamicki grafovi 
\section{Анализа друштвених мрежа}

Једна од главних мрежа података су друштвене мреже. Друштвена мрежа састоји се од релација између различитих друштвених ентитета. Иако као друштвени ентитет најчешће видимо особе, друштвени ентитети такође су и групе људи, организације и било које друге групације које су повезане неким односима или интеракцијама између њих. Природно се овакве мреже моделују графовима, где су чворови друштвени ентитети, а гране представљају односе иземеђу њих~\cite{tabassumSocialNetworkAnalysis2018}. Друштвене мреже карактерише велики број честих промена па су динамички графови идеалне структуре за њихово моделовање~\cite{APPLICATIONS_OF_DYNAMIC_GRAPH}.

Над оваквим графовима можемо вршити разне анализе и испитивати односе међу ентитетима, пратити промене у тим односима и на основу тога доносити закључке о друштву и ентитетима у њему. На пример можемо испитивати који су ентитети најпознатији, најмање познати, који ентитеети повезују различите групе и слично. Ове особине лако се преносе на особине чворова у графу, рецимо уколико повезаност чворова представља релацију "познаје" и нека је та релација симетрична, тада је чвор највећег степена уједно и најпознатији ентитет у друштву. 

Једно својство специфично за друштвене мреже је тенденција креирања заједница (енгл. \textit{communitygraph}). У контексту графова заједнице видимо као подграфове у којима су чворови густо повезани, а шворови из различитих подграфова ретко повезани. На слици \ref{fig:community_graph} можемо видети пример друштва са дванаест ентитета и три заједнице~\cite{tabassumSocialNetworkAnalysis2018}.


\begin{figure}[ht]
\centering
\begingroup
\shorthandoff{"}
\digraph[scale=0.6]{communitygraph}{
    node [shape=circle];
    edge [dir=none];

    subgraph cluster_A {
        style=filled;
        color=lightgrey;

        A1 -> A3;
        A2 -> A3;
        A3 -> A4;
        A4 -> A1;
        A1 -> A3;
    }

    subgraph cluster_B {
        style=filled;
        color=lightblue;
        B1 -> B2;
        B2 -> B3;
        B3 -> B4;
        B1 -> B4;
        B2 -> B4;
    }

    subgraph cluster_C {
        style=filled;
        color=lightyellow;
        C1 -> C2;
        C2 -> C3;
        C4 -> C1;
        C2 -> C4;
    } 

    A4 -> B1;
    B4 -> C2 [constraint=false];
    A2 -> C3 [constraint=false];
}
\endgroup
\caption{Друштво са 12 ентитета и 3 заједнице}
\label{fig:community_graph}
\end{figure}

\section{Интернет и Google PageRank}

Развојем интернета број страница и хиперлинкова између њих је растао великом брзином. Структура која се користи за чување и анализу ових података зове се Веб граф (енгл. \textit{The Web graph}). Веб граф представља усмерени граф чији су чворови интернет странице, а два чвора $u$ и $v$ су повезана граном уколико страница $u$ садржи хиперлинк на страницу $v$. Веб граф је огроман објекат, тренутно садржи преко 3 милијарде чворова повезаних са преко 50 милијарди грана~\cite{boldiWebgraphFrameworkCompression2004}. Како се у времену интернета број страница и хиперлинкови између њих мењају свакодневно динамички графови представљају погодан избор за ефикасно чување и анализу ових података. 

\subsection{Google PageRank}

PageRank је алгоритам за анализу линкова који се заснива на концепту централности сопствених вектора (\textit{eigenvector centrality}). Овај алгоритам користи интернет претраживач Google како би рангирао веб странице на основу вредности информација које оне носе, тако да се највредније странице приказују на самом врху резултата претраге. Основна идеја алгоритма јесте да се информације на вебу могу рангирати према популарности линкова. Што више веб страница показује на одређену страницу, то је она популарнија. Међутим, у овом процесу вредновања није релевантан само број линкова (односно степен чвора), већ и значај самих страница које упућују на њу. Стога PageRank мери релативни значај скупа веб страница не само на основу квантитета, већ пре свега на основу квалитета њихових линкова~\cite{boldiWebgraphFrameworkCompression2004}. Алгоритам PageRank један је од најпопуларнијих алгоритама на свету стога је вредно поменути га, али за потребе овог рада нећемо га детаљније објашњавати.

% Dynamic connectivity incremental
\section{Ширење заразних болести}

Болести као што су прехлада, грип, велике богиње или SARS, преносе се ваздухом између две особе које у затвореном простору проведу довољно времена заједно. То значи да можемо претпоставити да ће већина инфекција настати на локацијама као што су канцеларије, тржни центри, јавни превоз и слично~\cite{toroczkaiProximityNetworksEpidemics2007}. Већина стандардних математичких модела за ширење заразних болести користи диференцијалне једначине засноване на претпоставкама о униформном мешању или насумичне моделе који описују процес контакта~\cite{eubankModellingDiseaseOutbreaks2004}. Ипак, праћењем кретања људи можемо генерисати контактну мрежу, односно бипартитни граф $G_{PL}$ који се састоји од две врсте чворова: особа ($P$) и локација ($L$)~\cite{toroczkaiProximityNetworksEpidemics2007, eubankModellingDiseaseOutbreaks2004}. 

Уколико је особа $p$ посетила локацију $l$, у овом бипартитном графу постоји грана између чворова $p$ и $l$. Чворови могу бити означени одговарајућим демографским или географским подацима, док се гранама могу доделити ознаке са временом доласка и одласка. Из оваквог бипартитног графа природно се добијају две пројекције повезивањем свих парова чворова који се у бипартитном графу налазе на растојању два. Као резултат добијају се два засебна графа: граф особа $G_P$ и граф локација $G_L$. У графу особа $G_P$, гране су означене временским интервалима током којих су се особе налазиле на истој локацији. Ради једноставности, често се посматра статичка пројекција $\hat{G}_P$, која се добија занемаривањем временских ознака~\cite{eubankModellingDiseaseOutbreaks2004}.

Како се људи крећу и мењају своје локације, тако се и контактна мрежа мења. Стога је динамички граф погодна структура за моделовање оваквих мрежа. С обзиром на то да се у оваквим мрежама чворови и гране могу само додавати, јер једном остварен контакт у историји преноса не може бити избрисан, прецизније можемо рећи да је контактна мрежа инкрементални динамички граф.

У оваквом систему можемо искористити алгоритам динамичке повезаности да бисмо пратили ширење инфекције. Користећи пројекцију графа $G_P$ можемо пратити како се инфекција шири кроз популацију, док коришћењем пројекције $G_L$ пратимо угроженост самих локација. Проверу да ли је инфекција потенцијално достигла одређену особу или локацију вршимо једноставним упитом \texttt{IsConnected}, испитујући да ли су иницијално заражена особа и посматрана особа повезане у графу $G_P$. Сличну проверу можемо вршити и за локације у графу $G_L$.


\section{Протеин протеин интеракције}

\cite{zhangConstructionDynamicProbabilistic2016}

% MST
\section{3D реконструкција слика}

% Suitor Matching
\section{Алокација радних оптерећења у cloud инфраструктури}
\section{Препоручивање}



% Primene dinamickih grafova